\documentclass[11pt]{amsart}
\ifdefined\XeTeXrevision\else\ifdefined\pdfoutput\pdfoutput=1\fi\fi

\usepackage[margin=1.2in]{geometry}
\usepackage{amsmath,amssymb,mathtools}
\usepackage{booktabs}
\usepackage{array}
\usepackage{xcolor}
\definecolor{linkblue}{RGB}{25,60,120}
\usepackage[colorlinks=true,linkcolor=linkblue,citecolor=linkblue,urlcolor=linkblue]{hyperref}
\hypersetup{
  pdftitle={Peaceable queens on even tori: an exact parity formula},
  pdfauthor={Jan Philipp Harries},
  pdfsubject={Peaceable queens on toroidal chessboards},
  pdfkeywords={peaceable queens, toroidal chessboard, exact formula, moment relaxation, exact rational certificate, computer-assisted proof, OEIS A279405}
}

\newtheorem{theorem}{Theorem}
\newtheorem{lemma}[theorem]{Lemma}
\newtheorem{proposition}[theorem]{Proposition}
\newtheorem{corollary}[theorem]{Corollary}
\theoremstyle{definition}
\newtheorem{definition}[theorem]{Definition}
\theoremstyle{remark}
\newtheorem{remark}[theorem]{Remark}

\newcommand{\Z}{\mathbb{Z}}
\newcommand{\Q}{\mathbb{Q}}
\newcommand{\Zn}{\Z/n\Z}
\newcommand{\Bl}{\mathrm{B}}
\newcommand{\Wh}{\mathrm{W}}
\newcommand{\OPT}{\operatorname{OPT}_{42}}
\newcommand{\bs}{\discretionary{}{}{}}
\newcommand{\hs}[1]{\texttt{#1}}

\title[Peaceable queens on even tori]{Peaceable queens on even tori:\\an exact parity formula}

\author{Jan Philipp Harries}
\thanks{ellamind GmbH, Bremen, Germany. \textit{E-mail address:} \href{mailto:jan@ellamind.com}{jan@ellamind.com}}

\date{July 27, 2026}

\subjclass[2020]{Primary 05B99; Secondary 05C69, 68R05, 05C30, 90C05}
\keywords{peaceable queens, toroidal chessboard, parity construction, exact formula, moment relaxation, computer-assisted proof, OEIS A279405}

\raggedbottom

\begin{document}

\begin{abstract}
Let $t(n)$ be the largest $m$ for which $m$ black and $m$ white queens can be placed on the $n\times n$ toroidal chessboard with no black--white attack. We prove that, for every positive integer $q$,
\[
t(2q)=
\max_{(r_0,r_1,c_0,c_1)\in\{0,\ldots,q\}^4}
\min\!\left\{r_0c_1+r_1c_0,\,
(q-r_0)(q-c_0)+(q-r_1)(q-c_1)\right\}.
\]
The maximizing construction separates the two square parities and extends the
plaid idea of Clinch, Drescher, Huynh and Saffidine.  The upper bound combines
an exact moment relaxation, rational branch certificates, a hand-checkable
algebraic finish for $q\ge130$, and a rerunnable exact integer cut-envelope
computation for the finite range.  In particular,
\[
\lim_{\substack{n\to\infty\\2\mid n}}\frac{t(n)}{n^2}
=\frac{2-\sqrt3}{2}=0.1339745\ldots .
\]
We also prove the odd values $t(15)=20$ and $t(17)=28$ by exhaustive finite
searches with separately written audits.
\end{abstract}

\maketitle

\section{Introduction}\label{sec:intro}

A peaceable placement consists of equally many black and white queens, with no
queen attacking a queen of the other colour; attacks within an army are
allowed.  The problem goes back to Ainley \cite{Ainley77}.  On the ordinary
board its asymptotic density remains open \cite{CDHS25}; on the toroidal board
the rows, columns and diagonals wrap around, and the problem becomes
translation invariant.  Its maxima form OEIS sequence A279405
\cite{OEIS279405}.

Clinch, Drescher, Huynh and Saffidine \cite{CDHS25,CDHS24} introduced
\emph{plaid} constructions on even tori.  They obtained asymptotic lower
density $0.1339$ and upper density $0.1402$, while for odd orders they proved
the upper bound $t(n)\le n^2/8$.  The following theorem determines the even
problem exactly.

\begin{definition}\label{def:H}
For a positive integer $q$, put
\begin{equation}\label{eq:Hdef}
H(2q):=
\max_{(r_0,r_1,c_0,c_1)\in\{0,\ldots,q\}^4}
\min\!\left\{r_0c_1+r_1c_0,\,
(q-r_0)(q-c_0)+(q-r_1)(q-c_1)\right\}.
\end{equation}
\end{definition}

\begin{theorem}\label{thm:main}
For every positive integer $q$,
\[
t(2q)=H(2q).
\]
\end{theorem}

\begin{corollary}\label{cor:asym}
Let $\kappa=4-2\sqrt3$ and $\gamma=1-1/\sqrt3$.  For every positive
integer $q$,
\[
\kappa q^2-\gamma q-\tfrac14\le H(2q)\le\kappa q^2.
\]
Consequently, for even $n$,
\[
t(n)=\tfrac{2-\sqrt3}{2}\,n^2+O(n),
\qquad
\lim_{\substack{n\to\infty\\2\mid n}}\frac{t(n)}{n^2}
=\frac{2-\sqrt3}{2}=0.1339745\ldots,
\]
and $t(n)\le\lfloor(2-\sqrt3)n^2/2\rfloor$.
\end{corollary}

\begin{corollary}\label{cor:values}
Every even-indexed term $t(2q)$ of A279405 is computable from \eqref{eq:Hdef}; for example,
\[
t(16)=32,\quad t(18)=42,\quad t(20)=51,\quad t(40)=210 .
\]
\end{corollary}

\begin{theorem}\label{thm:oddmain}
$t(15)=20$ and $t(17)=28$.
\end{theorem}

Theorem~\ref{thm:main} implies that there always exists an optimal
\emph{parity-separated} line colouring: the diagonal colours reserve one
square parity for each army, leaving the four row and column parity counts in
\eqref{eq:Hdef}.  This family contains the classical plaids but is strictly
larger; for example it gives $H(20)=51$, whereas direct optimization of the
classical two-parameter subfamily gives $50$.

For the upper bound we first replace queens by a colouring of the four line
families.  On an even torus, conditioning on row and column parities yields a
$42$-equation moment relaxation.  Two exact rational branch certificates
localize every possible counterexample to a small neighbourhood of the
parity-separated profile.  Two support inequalities then reduce that
neighbourhood to a short algebraic argument for $q\ge130$.  The remaining
orders are covered by a rerunnable exact integer cut-envelope computation,
apart from five small orders settled by direct exhaustive search.
Section~\ref{sec:odd} gives a proof-complete account of the searches for
Theorem~\ref{thm:oddmain}.
Appendix~\ref{app:verify} describes the released certificates, audits and
reproduction commands.  All proof decisions use integer or rational
arithmetic.

\section{The formula and the lower bound}\label{sec:formula}

\subsection{Line colourings}

Throughout, $n\ge2$, all coordinates lie in $\Zn=\{0,1,\dots,n-1\}$, and the board is $G=(\Zn)^2$. Through a cell $(r,c)$ pass four \emph{lines}: row $r$, column $c$, diagonal $d=r-c$ and antidiagonal $a=r+c$, each family indexed by $\Zn$. Two queens attack each other exactly when they share a line. For $R,C,D,A\subseteq\Zn$ put
\[
\Bl(R,C,D,A)=\{(r,c)\in G:\ r\in R,\ c\in C,\ r-c\in D,\ r+c\in A\},
\]
and $\Wh(R,C,D,A)=\Bl(R^{c},C^{c},D^{c},A^{c})$, complements taken family by family.

\begin{lemma}[Line-colouring equivalence]\label{lem:lines}
Let $m\ge1$. A peaceable placement of $m$ black and $m$ white queens on the $n\times n$ torus exists if and only if there are $R,C,D,A\subseteq\Zn$ with $|\Bl|\ge m$ and $|\Wh|\ge m$.
\end{lemma}

\begin{proof}
($\Rightarrow$) Colour every line carrying a black queen black and every other line white; no line carries queens of both colours, else those two queens attack. Every black queen then lies on four black lines and every white queen on four white lines, and the queens of one colour occupy distinct cells.

($\Leftarrow$) Place black queens on any $m$ cells of $\Bl$ and white queens on any $m$ cells of $\Wh$. The two sets are disjoint, and a black--white attack would need a line lying simultaneously in a black set and in its complement.
\end{proof}

Consequently $t(n)=\max_{R,C,D,A}\min(|\Bl|,|\Wh|)$, and the queens have disappeared. Complementing all four sets exchanges $\Bl$ and $\Wh$, so $\min(|\Bl|,|\Wh|)$ is invariant under complementation; we use this freely.

\begin{lemma}[Incidence]\label{lem:inc}
Let $\ell,\ell'$ be lines of two distinct families.  Unless they form a
diagonal--antidiagonal pair, they meet in exactly one cell.  A diagonal $d$
and an antidiagonal $a$ meet in the fibre
\[
\varphi^{-1}(d,a),\qquad \varphi(r,c)=(r-c,r+c).
\]
Every point in this fibre satisfies $2r=a+d$ and $2c=a-d$; these doubled
congruences are equivalent to the original pair only when $n$ is odd.  The
fibre has one cell for odd $n$.  For even $n$ it has two cells, differing by
$(n/2,n/2)$, when $d\equiv a\pmod2$, and is empty otherwise.
\end{lemma}

\begin{proof}
The first statement is immediate (a row $r$ and a column $c$ meet in $(r,c)$; a row $r$ and a diagonal $d$ in $(r,r-d)$; and so on). For the second, consider $\varphi(r,c)=(r-c,r+c)$ on $G$. Its kernel is $\{(r,c):r=c,\ 2r=0\}$, trivial for odd $n$ and of order $2$ for even $n$, generated by $(\frac n2,\frac n2)$. For even $n$ every nonempty fibre thus has two elements and the image has $n^2/2$ elements; every image point satisfies $d\equiv a\pmod 2$, and there are exactly $n^2/2$ such pairs, so the image is precisely that set.
\end{proof}

\begin{lemma}[Budget]\label{lem:budget}
If $|\Bl|\ge m$ and $|\Wh|\ge m$ for some $(R,C,D,A)$, then the same holds
for a configuration satisfying
\[
|R|+|C|+|D|+|A|\le2n.
\]
\end{lemma}

\begin{proof}
Complementing all four sets exchanges $\Bl$ and $\Wh$ and replaces the sum
of their sizes by $4n$ minus that sum.
\end{proof}

For two line families $X,Y$, let $I(X,Y)$ denote the number of cells whose
lines in both families are black.

\begin{lemma}[Triple identity]\label{lem:triple}
Choose three line families with black sets of sizes $x,y,z$.  If $U$ and
$U'$ count the cells whose three lines are all black and all white,
respectively, then
\begin{equation}\label{eq:triple}
U+U'=n^2-n(x+y+z)+I(X,Y)+I(X,Z)+I(Y,Z).
\end{equation}
In particular, $|\Bl|,|\Wh|\ge m$ implies that the right-hand side is at
least $2m$.
\end{lemma}

\begin{proof}
For the three black-line indicators $u,v,w$ on the cells,
\[
uvw+(1-u)(1-v)(1-w)=1-u-v-w+uv+uw+vw.
\]
Summing gives \eqref{eq:triple}; moreover, $\Bl\subseteq U$ and
$\Wh\subseteq U'$.
\end{proof}

\subsection{The parity-separated construction}

Now let $n=2q$.  The diagonal and antidiagonal labels through $(r,c)$ have
the same parity, namely that of $r+c$.  Colouring both diagonal families
black on one label parity and white on the other therefore reserves one
square parity for each army.  Only the row and column choices remain.

Split each of the four families by label parity: for $i\in\{0,1\}$ let $r_i=|\{\rho\in R:\rho\equiv i\}|$, and likewise $c_i,d_i,a_i$; each parity class of $\Z/2q\Z$ has exactly $q$ labels, so all eight counts lie in $\{0,\dots,q\}$.

\begin{proposition}[Parity separation]\label{prop:lower}
Let $q\ge1$ and $0\le r_0,r_1,c_0,c_1\le q$. Choose $R\subseteq\Z/2q\Z$ consisting of any $r_0$ even and any $r_1$ odd labels, $C$ of any $c_0$ even and $c_1$ odd labels, and set $D=A=\{\text{odd labels}\}$. Then
\[
|\Bl|=r_0c_1+r_1c_0,
\qquad
|\Wh|=(q-r_0)(q-c_0)+(q-r_1)(q-c_1).
\]
Consequently $t(2q)\ge H(2q)$.
\end{proposition}

\begin{proof}
A cell $(r,c)$ is black-homogeneous iff $r\in R$, $c\in C$ and both $r-c$ and $r+c$ are odd, i.e.\ iff $r\in R$, $c\in C$ and $r,c$ have opposite parities; there are $r_0c_1+r_1c_0$ such cells. It is white-homogeneous iff $r\notin R$, $c\notin C$ and $r+c$ is even, i.e.\ $r,c$ of equal parity; there are $(q-r_0)(q-c_0)+(q-r_1)(q-c_1)$ such cells. Now apply Lemma~\ref{lem:lines} with $m=\min(|\Bl|,|\Wh|)$ and maximize over the four counts.
\end{proof}

\begin{remark}\label{rem:plaid16}
The construction extends the plaid mechanism of \cite[\S3]{CDHS24}.
Up to exchanging parity classes, the classical two-parameter subfamily has
profiles $(u,q,v,q)$ and hence capacity
\[
\min\{q(u+v),(q-u)(q-v)\},\qquad 0\le u,v\le q.
\]
The four counts in Proposition~\ref{prop:lower} are independent, so its
family is larger.  Maximizing the displayed two-variable expression gives
$50$ at $n=20$ and $72$ at $n=24$, whereas $H(20)=51$ and $H(24)=74$.
At $n=16$, the profile $(0,6,6,0)$ gives $36$ black-homogeneous and $32$
white-homogeneous cells.
\end{remark}

\subsection{Values and size of \texorpdfstring{$H$}{H}}

\begin{table}[ht]
\centering\small
\begin{tabular}{@{}l*{10}{r}@{}}
\toprule
$n$      & 2 & 4 & 6 & 8 & 10 & 12 & 14 & 16 & 18 & 20\\
$H(n)$   & 0 & 2 & 4 & 8 & 12 & 18 & 25 & 32 & 42 & 51\\
\midrule
$n$      & 22 & 24 & 26 & 28 & 30 & 32 & 34 & 36 & 38 & 40\\
$H(n)$   & 64 & 74 & 90 & 101 & 120 & 133 & 153 & 170 & 190 & 210\\
\bottomrule
\end{tabular}
\medskip
\caption{Exact even values through order 40.  The entries through order 14
agree with A279405, accessed July 27, 2026.}\label{tab:Hvals}
\end{table}

\begin{lemma}[Continuous upper bound]\label{lem:Hupper}
Let $q>0$ be real and let $r_0,r_1,c_0,c_1\in[0,q]$ be real. Put
\[
X=r_0c_1+r_1c_0,\qquad Y=(q-r_0)(q-c_0)+(q-r_1)(q-c_1).
\]
Then $\min(X,Y)\le\kappa q^2$, where $\kappa=4-2\sqrt3$. In particular $H(2q)\le\lfloor\kappa q^2\rfloor=\bigl\lfloor\tfrac{2-\sqrt3}{2}(2q)^2\bigr\rfloor$.
\end{lemma}

\begin{proof}
Normalize: $\tilde r_i=r_i/q$, $\tilde c_i=c_i/q$, and set
\[
\rho=\tilde r_0+\tilde r_1,\quad \chi=\tilde c_0+\tilde c_1,\quad
\xi=1-\rho,\quad\eta=1-\chi,\quad
z=(\tilde r_0-\tilde r_1)(\tilde c_0-\tilde c_1).
\]
Expanding,
\begin{equation}\label{eq:XYid}
\frac{X-Y}{q^2}=\rho+\chi-2-z,
\qquad
\frac{X+Y}{q^2}=2-\rho-\chi+\rho\chi=1+\xi\eta ,
\end{equation}
and $|z|\le\rho\chi$ because $|\tilde r_0-\tilde r_1|\le\tilde r_0+\tilde r_1=\rho$ and likewise for the columns.

If $\xi\eta\le0$ then $\min(X,Y)\le\frac{X+Y}{2}\le\frac{q^2}{2}<\kappa q^2$.

The map $(r_0,r_1,c_0,c_1)\mapsto(q-r_0,q-r_1,q-c_1,q-c_0)$ preserves the domain, exchanges $X$ and $Y$, and replaces $(\xi,\eta)$ by $(-\xi,-\eta)$. So we may assume $\xi,\eta\ge0$.

Suppose first $X\ge Y$. Then $z\le\rho+\chi-2$ by \eqref{eq:XYid}, and $z\ge-\rho\chi$, so $2-\rho-\chi\le\rho\chi$, i.e.\ $\xi+\eta\le1-\xi-\eta+\xi\eta$, i.e.
\begin{equation}\label{eq:xieta}
2\xi+2\eta-\xi\eta\le1 .
\end{equation}
Put $w=\sqrt{\xi\eta}\in[0,1]$. By AM--GM, $4w-w^2\le2\xi+2\eta-\xi\eta\le1$, so $w^2-4w+1\ge0$ and hence $w\le2-\sqrt3$. By \eqref{eq:XYid},
\[
\min(X,Y)\le\frac{X+Y}{2}=q^2\,\frac{1+w^2}{2}\le q^2\,\frac{1+(2-\sqrt3)^2}{2}=(4-2\sqrt3)q^2 .
\]

Now suppose $Y\ge X$, so $\min(X,Y)=X$, and assume for contradiction $X>\kappa q^2$. Then $X+Y\ge2X>2\kappa q^2$, so by \eqref{eq:XYid} $\xi\eta>2\kappa-1=(2-\sqrt3)^2$ and $w>2-\sqrt3$. On the other hand $X/q^2=\tilde r_0\tilde c_1+\tilde r_1\tilde c_0$ is bilinear in $(\tilde r_0,\tilde r_1)$ and $(\tilde c_0,\tilde c_1)$ subject to the fixed sums $\rho,\chi$, hence is at most $\rho\chi$ (attained at a vertex), and $\rho\chi=(1-\xi)(1-\eta)\le(1-w)^2$ by AM--GM. Therefore
\[
\kappa<\frac{X}{q^2}\le(1-w)^2<\bigl(1-(2-\sqrt3)\bigr)^2=(\sqrt3-1)^2=4-2\sqrt3=\kappa,
\]
a contradiction. Finally, $H(2q)$ is an integer maximum of $\min(X,Y)$ over a subset of the domain, so $H(2q)\le\lfloor\kappa q^2\rfloor$.
\end{proof}

\begin{lemma}\label{lem:Hlower}
Let $s=\sqrt3-1$, so $\kappa=s^2=4-2\sqrt3$, and let $\gamma=1-1/\sqrt3=s/(1+s)$. Then for every $q\ge1$,
\[
H(2q)\ \ge\ \kappa q^2-\gamma q-\tfrac14 .
\]
\end{lemma}

\begin{proof}
For an integer $0\le P\le2q$ take the profile $(r_0,r_1,c_0,c_1)=\bigl(0,\lfloor P/2\rfloor,\lceil P/2\rceil,0\bigr)$. Then
\[
X=\Bigl\lfloor\frac{P}{2}\Bigr\rfloor\Bigl\lceil\frac{P}{2}\Bigr\rceil=\Bigl\lfloor\frac{P^2}{4}\Bigr\rfloor,
\qquad
Y=q\Bigl(q-\Bigl\lceil\frac P2\Bigr\rceil\Bigr)+\Bigl(q-\Bigl\lfloor\frac P2\Bigr\rfloor\Bigr)q=2q^2-qP,
\]
so $H(2q)\ge\min(\lfloor P^2/4\rfloor,\ 2q^2-qP)$. The two branches cross where $P^2/4=2q^2-qP$, i.e.\ at $P=2sq$, with common value $s^2q^2=\kappa q^2$.

Write $2sq=m+\theta$ with $m=\lfloor2sq\rfloor$ and $\theta\in[0,1)$. Taking $P=m\le2sq$ gives $X\ge P^2/4-\frac14=\kappa q^2-sq\theta+\theta^2/4-\frac14\ge\kappa q^2-sq\theta-\frac14$ and $Y=2q^2-qm=\kappa q^2+q\theta\ge\kappa q^2$, so the value is at least $\kappa q^2-sq\theta-\frac14$. Taking $P=m+1\ge2sq$ gives $X\ge P^2/4-\frac14\ge\kappa q^2-\frac14$ and $Y=\kappa q^2-q(1-\theta)$, so the value is at least $\kappa q^2-q(1-\theta)-\frac14$. Hence
\[
H(2q)\ \ge\ \kappa q^2-q\min(s\theta,\,1-\theta)-\tfrac14\ \ge\ \kappa q^2-q\max_{0\le\theta<1}\min(s\theta,1-\theta)-\tfrac14 ,
\]
and the inner maximum is attained at $s\theta=1-\theta$, with value $s/(1+s)=1-1/\sqrt3=\gamma$.
\end{proof}

Lemmas~\ref{lem:Hupper} and~\ref{lem:Hlower} give Corollary~\ref{cor:asym} once Theorem~\ref{thm:main} is proved. Numerically $\gamma=0.42264\ldots$, so $H(2q)$ deviates from $\kappa q^2$ by less than $\tfrac12 q+\tfrac14$.

\section{The upper bound: relaxation and canonical chamber}\label{sec:relax}

From here to the end of Section~\ref{sec:local}, $n=2q$ and we prove $t(2q)\le H(2q)$ for $q\ge130$. Fix a configuration $(R,C,D,A)$ with black count $B=|\Bl|$ and white count $W=|\Wh|$, and let
\[
(r_0,r_1,c_0,c_1,d_0,d_1,a_0,a_1)\in\{0,\dots,q\}^8
\]
be its eight parity-split outer counts, with normalization $x=(r_0,\dots,a_1)/q\in[0,1]^8$.

\subsection{The \texorpdfstring{$42$}{42}-equation moment relaxation}

Partition the board into four \emph{row/column parity blocks} $G_{ij}=\{(r,c):r\equiv i,\ c\equiv j\pmod2\}$, each of $q^2$ cells. Inside $G_{ij}$ every cell has square parity $e=i\oplus j$, so its diagonal and antidiagonal labels both have parity $e$; and each of the four line families restricted to the relevant parity class has exactly $q$ labels.

For each block introduce $16$ nonnegative variables
\[
p_{ij}(\varepsilon_R,\varepsilon_C,\varepsilon_D,\varepsilon_A),\qquad(\varepsilon_R,\varepsilon_C,\varepsilon_D,\varepsilon_A)\in\{0,1\}^4,
\]
intended as $q^{-2}$ times the number of cells of $G_{ij}$ whose row, column, diagonal and antidiagonal are black ($1$) or white ($0$) as indicated. Write $E_{ij}[\cdot]$ for the corresponding expectation. A genuine configuration satisfies the following $42$ equations, in which $e=i\oplus j$.

\begin{lemma}[Soundness of the moment system]\label{lem:moments}
For every configuration and every block $(i,j)$,
\begin{gather}
\textstyle\sum p_{ij}=1;\qquad
E_{ij}[\varepsilon_R]=\tfrac{r_i}{q},\quad E_{ij}[\varepsilon_C]=\tfrac{c_j}{q},\quad E_{ij}[\varepsilon_D]=\tfrac{d_e}{q},\quad E_{ij}[\varepsilon_A]=\tfrac{a_e}{q};
\label{eq:mom1}\\[2pt]
E_{ij}[\varepsilon_R\varepsilon_C]=\tfrac{r_ic_j}{q^2},\quad
E_{ij}[\varepsilon_R\varepsilon_D]=\tfrac{r_id_e}{q^2},\quad
E_{ij}[\varepsilon_R\varepsilon_A]=\tfrac{r_ia_e}{q^2},
\label{eq:mom2}\\[2pt]
E_{ij}[\varepsilon_C\varepsilon_D]=\tfrac{c_jd_e}{q^2},\quad
E_{ij}[\varepsilon_C\varepsilon_A]=\tfrac{c_ja_e}{q^2};
\notag
\end{gather}
and, for each square parity $e\in\{0,1\}$, the \emph{aggregate} equation
\begin{equation}\label{eq:mom3}
\sum_{i\oplus j=e}E_{ij}[\varepsilon_D\varepsilon_A]=\frac{2\,d_ea_e}{q^2} .
\end{equation}
These are $4\cdot10+2=42$ equations, whose right-hand sides involve exactly $22$ distinct products of outer counts. Moreover
\begin{equation}\label{eq:BW}
\frac{B}{q^2}=\sum_{i,j}p_{ij}(1,1,1,1),\qquad \frac{W}{q^2}=\sum_{i,j}p_{ij}(0,0,0,0).
\end{equation}
\end{lemma}

\begin{proof}
The normalization and the four marginals in \eqref{eq:mom1} hold because within $G_{ij}$ each of the four families' relevant parity classes is swept uniformly: for instance the map $G_{ij}\to(\Z_q)$, $(r,c)\mapsto$ (index of $r$ within its parity class), has all fibres of size $q$.

For \eqref{eq:mom2} we need each pair of families other than $\{D,A\}$ to fibre uniformly over $G_{ij}$. Writing $r=2r'+i$, $c=2c'+j$ with $r',c'\in\Z_q$, the pair $(\text{row},\text{column})$ corresponds to $(r',c')$, the pair (row, diagonal) to $(r',r'-c')$, (row, antidiagonal) to $(r',r'+c')$, (column, diagonal) to $(c',r'-c')$ and (column, antidiagonal) to $(c',r'+c')$. Each of these is, up to a block-dependent translation, an invertible linear map $\Z_q^2\to\Z_q^2$ of determinant $\pm1$, hence a bijection; so all $q^2$ label pairs occur exactly once and the corresponding moment factorizes into the product of the two black-line densities.

The pair $\{D,A\}$ does not fibre uniformly within a block: $(r',c')\mapsto(r'-c',r'+c')$ has determinant $2$. But by Lemma~\ref{lem:inc} each pair $(d,a)$ of labels of equal parity $e$ meets in exactly two cells, and these two cells $(r,c)$, $(r+q,c+q)$ lie in the two blocks with $i\oplus j=e$ (in the same block iff $q$ is even, in different blocks otherwise --- in either case in the union). Summing over the two blocks therefore counts each equal-parity label pair exactly twice, giving \eqref{eq:mom3}. Finally \eqref{eq:BW} is the definition.
\end{proof}

\begin{remark}
The failure of the $\{D,A\}$ pair to factorize \emph{per block} is not cosmetic: for $4\mid n$ the two lifts of an equal-parity $(d,a)$ pair land in the \emph{same} block, so per-block independence of the diagonal families is false and $2$-adic information survives parity conditioning. Replacing \eqref{eq:mom3} by four per-block equations would be unsound. This is the reason the relaxation has $42$ and not $44$ rows.
\end{remark}

\begin{definition}\label{def:OPT}
For $x\in[0,1]^8$ let $\OPT(x)$ be the maximum of $t$ subject to: $p\ge0$; the $42$ equations of Lemma~\ref{lem:moments} with the right-hand sides evaluated at $x$ (that is, $r_i/q$ replaced by $x$-coordinates and each product $r_ic_j/q^2$ by the corresponding product of $x$-coordinates); and $t\le\sum_{ij}p_{ij}(1,1,1,1)$, $t\le\sum_{ij}p_{ij}(0,0,0,0)$.
\end{definition}

The maximum is well defined.  For each block and each $x$, four independent
Bernoulli indicators with the prescribed marginals satisfy all block
moments; summing the two relevant blocks also gives each aggregate
$\{D,A\}$ equation.  Thus the $p$-feasible set is nonempty.  It is compact,
and after eliminating $t$ the objective is the minimum of two continuous
linear functions of $p$, so the maximum is attained.

Lemma~\ref{lem:moments} immediately gives the key inequality
\begin{equation}\label{eq:opt}
\frac{\min(B,W)}{q^2}\ \le\ \OPT(x)
\end{equation}
for every configuration with normalized outer profile $x$.

\subsection{The canonical chamber}

\begin{lemma}[Chamber]\label{lem:chamber}
Every configuration can be transformed, without changing $\{B,W\}$ or the multiset of homogeneous cell counts, into one whose normalized profile $x$ satisfies
\begin{equation}\label{eq:chamber}
r_0\le r_1,\qquad c_0\le c_1,\qquad r_0+r_1\le c_0+c_1,\qquad d_0+d_1\le a_0+a_1,\qquad \sum_{i=1}^{8}x_i\le4 .
\end{equation}
\end{lemma}

\begin{proof}
Complementing all four line sets exchanges $B$ and $W$ and replaces $\sum x_i$ by $8-\sum x_i$, so the budget can be enforced. Translating the board by one unit in the row direction, $(r,c)\mapsto(r+1,c)$, exchanges the two row parity classes (and simultaneously the two diagonal and the two antidiagonal classes); translating by one unit in the column direction does the same for the columns. Composing the two if necessary, we may sort $r_0\le r_1$ and $c_0\le c_1$. Transposing, $(r,c)\mapsto(c,r)$, exchanges rows with columns and fixes the antidiagonals while negating the diagonals, hence orders the row and column totals. Finally $(r,c)\mapsto(r,-c)$ fixes the row parities, permutes the column labels within their parity classes, and exchanges the diagonal with the antidiagonal family, so it orders $d_0+d_1\le a_0+a_1$. Each of these is an affine bijection of the board mapping lines to lines, so $B$ and $W$ are preserved (or exchanged, in the complementation step).
\end{proof}

The four ordering inequalities in \eqref{eq:chamber}, apart from the budget
row, preserve the sorting achieved by the earlier steps.  Indeed, a unit
translation exchanges $(d_0,d_1)$ with $(d_1,d_0)$ \emph{and} $(a_0,a_1)$
with $(a_1,a_0)$ simultaneously, so it does not disturb the comparison of
the two family totals; the transpose does not affect either total.  We
therefore assume \eqref{eq:chamber} from now on.

\section{The two branch certificates and the escape comparison}\label{sec:cert}

\subsection{What the certificates are}

For a box $[l,u]\subseteq[0,1]^8$, replace each of the $22$ products $x_ix_j$ occurring in the right-hand sides of the moment equations by a fresh variable $y_{ij}$ constrained by the four McCormick envelope inequalities
\begin{equation}\label{eq:mccormick}
\begin{aligned}
-y_{ij}+l_ix_j+l_jx_i&\le l_il_j, & -y_{ij}+u_ix_j+u_jx_i&\le u_iu_j,\\
y_{ij}-u_ix_j-l_jx_i&\le-u_il_j, & y_{ij}-l_ix_j-u_jx_i&\le-l_iu_j,
\end{aligned}
\end{equation}
which are exactly the expansions of $(x_i-l_i)(x_j-l_j)\ge0$, $(u_i-x_i)(u_j-x_j)\ge0$, $(u_i-x_i)(x_j-l_j)\ge0$ and $(x_i-l_i)(u_j-x_j)\ge0$.  We also retain the redundant valid bounds
\[
l_il_j\le y_{ij}\le u_iu_j;
\]
these are used when maximizing the partially dualized Lagrangian.  Together
with the budget row, the four chamber rows of \eqref{eq:chamber}, and the two
rows $t\le B/q^2$, $t\le W/q^2$, this gives a linear programme in the
$8+22+64+1=95$ variables $(x,y,p,t)$ with $95$ non-bound inequality rows
and $42$ equality rows, whose optimum dominates $\OPT(x)$ for every $x$ in
the box.

\begin{lemma}[Dual-leaf soundness]\label{lem:dualleaf}
Consider the linear programme above on a box: maximize $t$ over
$z=(x,y,p,t)$ subject to the inequality rows $Az\le b$ (McCormick, budget,
chamber and the two objective rows), the $42$ moment equalities $Ez=e$, the
box bounds on $(x,y)$, and $p\ge0$.  Let $\mu\le0$ and $\nu$ be rational
multipliers for the inequality and the equality rows, and form the partial
Lagrangian
\[
L(z)=t+\mu^{\mathsf T}(Az-b)+\nu^{\mathsf T}(Ez-e).
\]
Suppose that, after collecting terms, the coefficient of $t$ in $L$ vanishes
and the coefficient vector $\omega$ of the atom variables satisfies
$\omega\le0$, so that $L(z)=\ell_0+\ell(x,y)+\omega^{\mathsf T}p$ with
$\ell_0=-\mu^{\mathsf T}b-\nu^{\mathsf T}e$ and $\ell$ linear.  (Here $b$,
$e$ are the constraint right-hand sides, local to this lemma.)  Then every
feasible $z$ satisfies
\[
t\ \le\ \ell_0+\max_{(x,y)\in\mathrm{box}}\ell(x,y),
\]
and the maximum is computed exactly by taking each coordinate at the box
endpoint matching the sign of its coefficient in $\ell$.
\end{lemma}

\begin{proof}
If $z$ is feasible then $\mu^{\mathsf T}(Az-b)\ge0$, because componentwise
both factors are nonpositive, and $\nu^{\mathsf T}(Ez-e)=0$; hence
$t\le L(z)$.  Since $\omega\le0$ and $p\ge0$, we get
$L(z)\le\ell_0+\ell(x,y)$.
Finally, a linear function on a product of intervals is maximized coordinate
by coordinate at interval endpoints.
\end{proof}

A \emph{branch certificate} is a rooted binary tree. A split node carries a coordinate index and an exact rational cut, and its two children are the parent box intersected with $x_d\le c$ and with $x_d\ge c$ (the children overlap on the cutting hyperplane, which is harmless and gives an exhaustive cover). A leaf is one of three kinds:

\begin{itemize}
\item a \emph{McCormick-dual leaf}, carrying an exact rational multiplier vector for the inequality rows (nonpositive) and an unrestricted one for the equality rows, together with the resulting exact rational bound $u$. The verifier reconstructs the partial Lagrangian, checks that the reduced cost of $t$ vanishes and that the reduced costs of all $64$ atom variables are $\le0$ (so that the inner maximum over $p\ge0$ is finite), evaluates the outer and product variables at their better box endpoints, recomputes the bound itself, and checks that it equals the stored $u$ and satisfies the required threshold --- that is, it verifies precisely the hypotheses of Lemma~\ref{lem:dualleaf}, whose conclusion certifies $u$ as an upper bound for the relaxation on the leaf box;
\item an \emph{exact empty leaf}, carrying one of five reasons (\texttt{rsort}, \texttt{csort}, \texttt{rcsum}, \texttt{dasum}, \texttt{budget}). Each is a purely interval-arithmetic obstruction to the chamber inequalities: e.g.\ \texttt{rsort} asserts $\max(l_{r_1},l_{r_0})>u_{r_1}$, which contradicts $r_0\le r_1$; and \texttt{budget} asserts that the sum of the four group minima forced by the preceding steps exceeds $4$. These are exact rational emptiness tests, never solver infeasibility claims;
\item a \emph{plaid leaf}, asserting that the box lies wholly inside one of the four rational plaid neighbourhoods. With $s=\sqrt3-1$, the four templates are obtained by choosing a square parity and a row parity; a box qualifies when, coordinate by coordinate, the template-$0$ coordinates have upper endpoint $\le\frac1{10}$, the template-$1$ coordinates have lower endpoint $\ge\frac{19}{20}$, and the two template-$s$ coordinates lie in $[\frac35,\frac78]$.
\end{itemize}

\begin{lemma}[Branch-certificate inference]\label{lem:branchsound}
Suppose every terminal box of a branch certificate is either empty, contained
in a declared exceptional region, or carries a valid dual upper bound $u$.
Then the root box is the union of the exceptional terminal boxes and boxes
on which the relaxation has value at most the corresponding $u$.
\end{lemma}

\begin{proof}
Induct upward from the leaves.  At a split $x_d\le c$ or $x_d\ge c$, the two
closed child boxes cover the parent.  The stated alternatives are therefore
preserved at every internal node and hence at the root.
\end{proof}

\begin{definition}\label{def:N}
The canonical rational plaid neighbourhood is
\begin{equation}\label{eq:N}
N:\quad 0\le r_0,c_0,d_1,a_1\le\tfrac1{10},\qquad \tfrac35\le r_1,c_1\le\tfrac78,\qquad \tfrac{19}{20}\le d_0,a_0\le1
\end{equation}
(in normalized coordinates). After the chamber restrictions \eqref{eq:chamber}, it is the only one of the four templates that survives.
\end{definition}

\subsection{The two certificates}

\begin{theorem}[Global localization]\label{thm:cert1}
For every $x\in[0,1]^8$ satisfying the chamber inequalities \eqref{eq:chamber} and lying outside $N$,
\[
\OPT(x)\ \le\ \frac{527}{1000}.
\]
\end{theorem}

\begin{theorem}[Local core]\label{thm:cert2}
For every $x\in N$ satisfying \eqref{eq:chamber}, either $\OPT(x)\le\frac{527}{1000}$, or the unnormalized counts satisfy the four \emph{aggregate core} inequalities
\begin{equation}\label{eq:core}
\frac{36}{25}q\ \le\ r_1+c_1\ \le\ \frac{37}{25}q,
\qquad
r_0+c_0\le\frac{9}{100}q,
\qquad
d_1+a_1\le\frac{9}{200}q,
\qquad
(q-d_0)+(q-a_0)\le\frac{1}{25}q .
\end{equation}
\end{theorem}

Both theorems are established by exact branch certificates, distributed as the ancillary files \texttt{even\_\bs{}c527\_\bs{}sym\_\bs{}mc.json.gz} and \texttt{even\_\bs{}local\_\bs{}core\_\bs{}mc.json.gz}.

\begin{remark}[Verification]\label{rem:certverify}
Two programs sharing no code check these trees, both using only exact rational arithmetic and neither trusting any stored number it can recompute. The delivered checker \texttt{verify\_\bs{}even\_\bs{}branch\_\bs{}certificate.py} rebuilds the moment system, every McCormick row, every dual, every split and every empty leaf, and ends in \texttt{ALL\_\bs{}EVEN\_\bs{}BRANCH\_\bs{}CERTIFICATES\_\bs{}OK}. The from-scratch verifier \texttt{verify\_\bs{}even\_\bs{}finite\_\bs{}independent.py} does not import it or the builders; it \emph{discovers} the $42$ rows by counting incidence fibres on real tori at $n=8,12,16,20$ (keeping a row only when all label tuples are attained with equal multiplicity, which is precisely the factorization statement, and finding by itself that $\{D,A\}$ is the unique ragged pair, thereby deriving the two aggregate rows rather than assuming them), re-validates them as exact identities on $200$ random colourings at $n=12,16,20,28$ ($8{,}400$ equation checks), re-derives the McCormick rows, recomputes every leaf bound from its own partial Lagrangian, and re-proves every emptiness claim by its own interval reasoning. It agrees with the delivered checker on every count in Table~\ref{tab:certs}, and rejects eight deliberately tampered copies of the first certificate. Appendix~\ref{app:verify} gives commands and digests.
\end{remark}

\begin{proof}[Proof of Theorems~\ref{thm:cert1} and~\ref{thm:cert2}]
The exact checks in Remark~\ref{rem:certverify} establish that every internal
node makes the declared exhaustive split and that every leaf has the stated
type.  By Lemma~\ref{lem:dualleaf}, the stored multipliers of each
McCormick-dual leaf certify its bound $u$.  In the first tree, the only exceptional leaves lie in \(N\); all other
leaves are empty or have relaxation value at most \(527/1000\).  In the
second, the only exceptional leaves satisfy \eqref{eq:core}; all other leaves
are empty or have relaxation value at most \(527/1000\).  Applying
Lemma~\ref{lem:branchsound} to the two root boxes proves the respective
claims.
\end{proof}

Their statistics are:
\begin{table}[ht]
\centering
\begin{tabular}{@{}lrr@{}}
\toprule
 & Theorem~\ref{thm:cert1} & Theorem~\ref{thm:cert2}\\
\midrule
root box & $[0,1]^8$ & $N$\\
nodes & $6{,}929$ & $4{,}423$\\
split nodes & $3{,}464$ & $2{,}211$\\
McCormick-dual leaves & $3{,}368$ & $1{,}639$\\
plaid / core leaves & $4$ & $563$\\
exact empty leaves & $93$ & $10$\\
maximum depth & $22$ & $29$\\
largest reconstructed leaf bound & $\tfrac{527}{1000}-1.92\cdot10^{-6}$ & $\tfrac{2329117}{4419584}=\tfrac{527}{1000}-8.53\cdot10^{-7}$\\
\bottomrule
\end{tabular}
\medskip
\caption{The two branch certificates. The $93$ empty leaves of the first tree split as $1$ \texttt{rsort}, $9$ \texttt{csort}, $19$ \texttt{rcsum}, $21$ \texttt{dasum}, $43$ \texttt{budget}; the $10$ of the second are all \texttt{rcsum}. The four plaid leaves of the first tree all carry the template $(0,s,0,s,1,0,1,0)$, i.e.\ the neighbourhood $N$; the $563$ core leaves of the second all lie wholly inside \eqref{eq:core}. The last row shows that both certificates are tight to about $2\cdot10^{-6}$.}\label{tab:certs}
\end{table}

\subsection{The escape comparison}

\begin{lemma}\label{lem:escape}
$\kappa-\frac{527}{1000}>\frac1{125}$ and $\gamma<\frac12$. Consequently, for every $q\ge64$,
\begin{equation}\label{eq:escape}
H(2q)-\frac{527}{1000}q^2\ >\ \frac{q^2}{125}-\frac q2-\frac14\ >\ 0 .
\end{equation}
\end{lemma}

\begin{proof}
The inequality $\kappa-\frac{527}{1000}>\frac1{125}$ reads $4-2\sqrt3>\frac{535}{1000}=\frac{107}{200}$, i.e.\ $693>400\sqrt3$, i.e.\ $480249>480000$. The inequality $\gamma<\frac12$ reads $\frac1{\sqrt3}>\frac12$, i.e.\ $2>\sqrt3$. The first inequality of \eqref{eq:escape} is then Lemma~\ref{lem:Hlower} combined with these two comparisons. For the second, note that $q\ge64$ gives
\[
\frac{q^2}{125}-\frac q2-\frac14\ \ge\ q\Bigl(\frac{64}{125}-\frac12\Bigr)-\frac14=\frac{3q}{250}-\frac14\ \ge\ \frac{192}{250}-\frac14\ >\ 0 . \qedhere
\]
\end{proof}

Combining \eqref{eq:opt} with Theorems~\ref{thm:cert1},~\ref{thm:cert2} and Lemma~\ref{lem:escape}:

\begin{corollary}\label{cor:escape}
Let $q\ge64$ and let a configuration in the canonical chamber have $\min(B,W)\ge H(2q)$. Then its normalized profile lies in $N$ and its unnormalized counts satisfy the aggregate core inequalities \eqref{eq:core}.
\end{corollary}

\section{The local finish}\label{sec:local}

Throughout this section $q\ge130$, the configuration lies in the canonical chamber, its normalized profile lies in the neighbourhood $N$ of \eqref{eq:N}, and the aggregate core \eqref{eq:core} holds. We rename the eight counts to the local coordinates
\begin{equation}\label{eq:locvars}
u=r_0,\quad R=r_1,\quad v=c_0,\quad C=c_1,\qquad
e=d_1,\quad g=q-d_0,\quad f=a_1,\quad h=q-a_0,
\end{equation}
so $u,v,e,f,g,h$ are small ``defects'' and $R,C$ are close to $sq$. All eight are integers in $[0,q]$. Put
\begin{equation}\label{eq:aggregates}
\begin{aligned}
\sigma&=R+C, & m&=u+v, & \delta&=e+f, & \zeta&=g+h,\\
\alpha&=2q-\sigma-m, & \beta&=\sigma-q, & Q&=2ef+2gh,\\
X&=uv+RC, & \multicolumn{4}{l}{Y=(q-u)(q-C)+(q-R)(q-v).}
\end{aligned}
\end{equation}
The core \eqref{eq:core} reads
\begin{equation}\label{eq:core2}
\tfrac{36}{25}q\le\sigma\le\tfrac{37}{25}q,\qquad m\le\tfrac{9}{100}q,\qquad \delta\le\tfrac{9}{200}q,\qquad \zeta\le\tfrac1{25}q ,
\end{equation}
and the chamber \eqref{eq:chamber} gives, in these coordinates,
\begin{equation}\label{eq:chamber2}
u+R\le v+C,\qquad e+h\le f+g,\qquad \sigma+m+\delta-\zeta\le2q .
\end{equation}

Three remarks on notation.  First, from here to the end of the section $R$
and $C$ denote the two integer counts of \eqref{eq:locvars}; the line sets of
Section~\ref{sec:formula} are not referred to again before
Section~\ref{sec:odd}.  Likewise the defect sums are written $\delta,\zeta$
to avoid collision with the antidiagonal counts $a_0,a_1$ and with the
right-hand side $b(x)$ of Section~\ref{sec:ladder}.  Second, $(X,Y)$ are
exactly the two capacities of the $H$-profile
$(r_0,r_1,c_0,c_1)=(u,R,C,v)$ in Definition~\ref{def:H}, in which the two
column parity counts are \emph{exchanged} relative to the configuration's own
quadruple $(u,R,v,C)$; the exchanged quadruple is again admissible in
\eqref{eq:Hdef}, so its capacities are $r_0c_1+r_1c_0=uv+RC=X$ and
$(q-r_0)(q-c_0)+(q-r_1)(q-c_1)=(q-u)(q-C)+(q-R)(q-v)=Y$, and
\begin{equation}\label{eq:XYH}
\min(X,Y)\ \le\ H(2q).
\end{equation}
Third, $\alpha$ and $\beta$ are the two ``slacks'' that the support cuts below will charge against.

\subsection{The two support cuts}

The proof uses exactly two members of an exactly certified library of $760$ support-dual cuts for the moment system (Section~\ref{sec:ladder} and Appendix~\ref{app:verify}); they are entries $609$ and $559$ of that library, transported to the coordinates \eqref{eq:locvars}. A \emph{support dual} is a pair $(\alpha_{\mathrm{cut}},\lambda)$ of exact rationals such that the corresponding combination of the $42$ moment equations dominates $\alpha_{\mathrm{cut}}B+(1-\alpha_{\mathrm{cut}})W$ columnwise on all $64$ atoms; the resulting polynomial inequality then holds for every genuine configuration, with no relaxation and no branching. Cut $609$ has $\alpha_{\mathrm{cut}}=1$ and cut $559$ has $\alpha_{\mathrm{cut}}=0$, so they bound $B$ and $W$ separately:

\begin{lemma}[Cuts $609$ and $559$]\label{lem:cuts}
For every configuration in the canonical chamber,
\begin{equation}\label{eq:FB}
B\ \le\ F_B:=X-\beta\zeta+Q,
\end{equation}
\begin{equation}\label{eq:FW}
W\ \le\ F_W:=Y-\alpha\delta+Q .
\end{equation}
In particular, writing $T:=\min(F_B,F_W)$, we have $\min(B,W)\le T$.
\end{lemma}

The two exact rational duals are shipped in \texttt{even\_\bs{}finite\_\bs{}support\_\bs{}cuts.json} and are byte-identical to entries $609$ and $559$ of the certified library \texttt{new\_\bs{}dual\_\bs{}cuts.json}; \texttt{verify\_\bs{}even\_\bs{}finite\_\bs{}algebra.py} recomputes both polynomial identities and the columnwise support property. Everything below is elementary algebra from \eqref{eq:core2}, \eqref{eq:chamber2}, \eqref{eq:FB} and \eqref{eq:FW}.

\subsection{Basic estimates and the defect dichotomy}

\begin{lemma}\label{lem:basic}
Under \eqref{eq:core2},
\begin{equation}\label{eq:albe}
\alpha\ \ge\ \tfrac{43}{100}q,\qquad \beta\ \ge\ \tfrac{11}{25}q,\qquad \alpha>\delta,\qquad\beta>\zeta .
\end{equation}
\end{lemma}

\begin{proof}
$\alpha=2q-\sigma-m\ge2q-\frac{37}{25}q-\frac{9}{100}q=\frac{200-148-9}{100}q=\frac{43}{100}q$, and $\beta=\sigma-q\ge\frac{36}{25}q-q=\frac{11}{25}q$. Finally $\delta\le\frac{9}{200}q<\frac{43}{100}q\le\alpha$ and $\zeta\le\frac1{25}q<\frac{11}{25}q\le\beta$.
\end{proof}

\begin{lemma}[Defect dichotomy]\label{lem:dichotomy}
Under \eqref{eq:core2},
\begin{equation}\label{eq:dicho}
2Q=4ef+4gh\ \le\ \delta^2+\zeta^2\ \le\ \alpha\delta+\beta\zeta,
\qquad\text{hence}\qquad
Q\le\max(\alpha\delta,\ \beta\zeta).
\end{equation}
\end{lemma}

\begin{proof}
$4ef\le(e+f)^2=\delta^2$ and $4gh\le(g+h)^2=\zeta^2$. By Lemma~\ref{lem:basic}, $\delta^2\le\alpha\delta$ and $\zeta^2\le\beta\zeta$. Finally $Q\le\frac12(\alpha\delta+\beta\zeta)\le\max(\alpha\delta,\beta\zeta)$.
\end{proof}

This dichotomy is the only case distinction the local proof needs. We treat the two orderings of $X$ and $Y$ separately.

\subsection{Case \texorpdfstring{$X\le Y$}{X <= Y}}\label{ssec:caseA}

\begin{proposition}[First defect case]\label{prop:caseA}
If $X\le Y$ then $T\le H(2q)$.
\end{proposition}

\begin{proof}
If $Q\le\beta\zeta$ then $F_B\le X$ by \eqref{eq:FB}, so $T\le X=\min(X,Y)\le H(2q)$ by \eqref{eq:XYH}.

Assume therefore
\begin{equation}\label{eq:A17}
Q>\beta\zeta .
\end{equation}
By Lemma~\ref{lem:dichotomy}, $Q\le\alpha\delta$.

\emph{Step 1: a bound on $F_B$.} Since $\zeta\le q/25$,
\begin{equation}\label{eq:A21a}
2gh\le\frac{\zeta^2}{2}\le\frac{q\zeta}{50}\le\beta\zeta ,
\end{equation}
the last step by Lemma~\ref{lem:basic}. Hence
\begin{equation}\label{eq:A21}
F_B=X-\beta\zeta+2ef+2gh\ \le\ X+2ef .
\end{equation}

\emph{Step 2: four switched profiles.} Consider the four $H$-profiles
\begin{equation}\label{eq:A18}
(r_0,r_1,c_0,c_1)=(u+x,\ R,\ C,\ v+y),
\qquad (x,y)\in\bigl\{(e,2f),\,(2e,f),\,(f,2e),\,(2f,e)\bigr\} .
\end{equation}
All coordinates lie in $[0,q]$: inside $N$ we have $u,v\le q/10$ and $e,f\le q/10$, so the largest modified coordinate is at most $q/10+2q/10=3q/10$. The black capacity of \eqref{eq:A18} is
\begin{equation}\label{eq:A20}
(u+x)(v+y)+RC=X+uy+vx+xy\ \ge\ X+xy\ =\ X+2ef ,
\end{equation}
since $u,v,x,y\ge0$ and each of the four choices in \eqref{eq:A18} has $xy=2ef$. Combining with \eqref{eq:A21}, each of the four profiles has black capacity at least $F_B$.

The white capacity of \eqref{eq:A18} is $(q-u-x)(q-C)+(q-R)(q-v-y)=Y-\bigl(x(q-C)+y(q-R)\bigr)$, so the \emph{white loss} is $x(q-C)+y(q-R)$. Put $L=2q-\sigma$. Over the four choices in \eqref{eq:A18} the values of $x$ sum to $3(e+f)=3\delta$ and so do the values of $y$; hence the average white loss is exactly
\begin{equation}\label{eq:A22}
\frac{3\delta}{4}\bigl((q-C)+(q-R)\bigr)=\frac34\,\delta L .
\end{equation}

\emph{Step 3: the averaged loss is affordable.} We claim
\begin{equation}\label{eq:A23}
Q\ \le\ \delta\Bigl(\frac L4-m\Bigr).
\end{equation}
First, by \eqref{eq:core2},
\begin{equation}\label{eq:A24}
\frac L4-m\ \ge\ \frac14\Bigl(2-\frac{37}{25}\Bigr)q-\frac9{100}q=\frac{13}{100}q-\frac{9}{100}q=\frac q{25}.
\end{equation}
Second, from \eqref{eq:A17}, Lemma~\ref{lem:basic} and \eqref{eq:A21a},
\[
2ef=Q-2gh>\beta\zeta-\frac{q\zeta}{50}\ \ge\ \frac{11}{25}q\zeta-\frac{q\zeta}{50}=\frac{21}{50}\,q\zeta .
\]
Since $2ef\le\delta^2/2$, this gives $\zeta<\dfrac{25\delta^2}{21q}$, and therefore
\begin{equation}\label{eq:A25}
Q\le\frac{\delta^2+\zeta^2}{2}\le\frac{\delta^2}{2}+\frac{q\zeta}{50}<\frac{\delta^2}{2}+\frac{q}{50}\cdot\frac{25\delta^2}{21q}=\frac{11}{21}\delta^2\le\frac{11}{21}\cdot\frac{9q}{200}\,\delta=\frac{33}{1400}\,\delta q<\frac{1}{25}\delta q ,
\end{equation}
using $\zeta^2/2\le q\zeta/50$ (valid as $\zeta\le q/25$) and $\delta\le9q/200$. Comparing \eqref{eq:A24} and \eqref{eq:A25} proves \eqref{eq:A23}.

\emph{Step 4: conclusion.} Since $\alpha=L-m$, inequality \eqref{eq:A23} is equivalent to
\begin{equation}\label{eq:A26}
\frac34\,\delta L\ \le\ \delta(L-m)-Q=\alpha\delta-Q .
\end{equation}
By \eqref{eq:A22} at least one of the four profiles has white loss at most $\frac34\delta L$, hence white capacity at least $Y-(\alpha\delta-Q)=F_W$. By \eqref{eq:A20} and \eqref{eq:A21} its black capacity is at least $F_B$. That profile therefore has $H$-value at least $\min(F_B,F_W)=T$, and so $T\le H(2q)$.
\end{proof}

\subsection{Case \texorpdfstring{$Y\le X$}{Y <= X}}\label{ssec:caseB}

\begin{proposition}[Second defect case]\label{prop:caseB}
If $Y\le X$ and $q\ge130$, then $T\le H(2q)$.
\end{proposition}

\begin{proof}
If $Q\le\alpha\delta$ then $F_W\le Y$ by \eqref{eq:FW}, so $T\le Y=\min(X,Y)\le H(2q)$ by \eqref{eq:XYH}.

Assume therefore
\begin{equation}\label{eq:B28}
Q>\alpha\delta .
\end{equation}
By Lemma~\ref{lem:dichotomy}, $Q\le\beta\zeta$.

\emph{Step 1: the integrality step.} We claim $gh>0$.  If $gh=0$, then
\eqref{eq:B28} forces $\delta>0$: otherwise Lemma~\ref{lem:basic} gives
$e=f=0$, and hence $Q=0$, contradicting $Q>\alpha\delta=0$.  Using
$\delta\le\frac{9}{200}q$ and Lemma~\ref{lem:basic}, we then have
\[
Q=2ef\le\frac{\delta^2}{2}\le\frac{9q}{400}\,\delta<\frac{43}{100}q\,\delta\le\alpha\delta ,
\]
contradicting \eqref{eq:B28}. Since $g$ and $h$ are integer counts, $g\ge1$ and $h\ge1$, so
\begin{equation}\label{eq:B29}
\zeta=g+h\ \ge\ 2 .
\end{equation}
This is the only point in the whole argument where integrality of the counts is used; it is also exactly the point at which the continuous relaxation is provably insufficient (Remark~\ref{rem:integrality}).

\emph{Step 2: averaging the two cuts.} By \eqref{eq:FB} and \eqref{eq:FW},
\begin{equation}\label{eq:B30}
T\ \le\ \frac{F_B+F_W}{2}=\frac{X+Y}{2}-\frac{\alpha\delta+\beta\zeta-2Q}{2}.
\end{equation}
By \eqref{eq:dicho},
\begin{equation}\label{eq:B31}
\alpha\delta+\beta\zeta-2Q\ \ge\ \alpha\delta+\beta\zeta-(\delta^2+\zeta^2)=\delta(\alpha-\delta)+\zeta(\beta-\zeta)\ \ge\ \zeta(\beta-\zeta),
\end{equation}
since $\delta(\alpha-\delta)\ge0$ by Lemma~\ref{lem:basic}. On the range $2\le\zeta\le q/25$ the function $\zeta\mapsto\zeta(\beta-\zeta)$ is increasing, because its derivative $\beta-2\zeta\ge\frac{11}{25}q-\frac{2}{25}q=\frac9{25}q>0$. Hence by \eqref{eq:B29},
\begin{equation}\label{eq:B32}
\frac{\alpha\delta+\beta\zeta-2Q}{2}\ \ge\ \frac{2(\beta-2)}{2}=\beta-2 .
\end{equation}

\emph{Step 3: the balanced-plaid baseline.} We claim
\begin{equation}\label{eq:B33}
\frac{X+Y}{2}\ \le\ \kappa q^2 .
\end{equation}
Normalize $(u,R,C,v)$ by $q$ and put, as in Lemma~\ref{lem:Hupper},
\[
\rho=\frac{u+R}{q},\quad \chi=\frac{v+C}{q},\quad \xi=1-\rho,\quad \eta=1-\chi,\quad
z=\Bigl(\frac uq-\frac Rq\Bigr)\Bigl(\frac Cq-\frac vq\Bigr).
\]
The neighbourhood $N$ gives $\rho\le\frac1{10}+\frac78<1$ and $\chi<1$, so $\xi,\eta>0$. By \eqref{eq:XYid},
\[
\frac{X-Y}{q^2}=\rho+\chi-2-z,\qquad |z|\le\rho\chi,\qquad \frac{X+Y}{q^2}=1+\xi\eta .
\]
The assumption $X\ge Y$ forces $-\rho\chi\le z\le\rho+\chi-2$, whence $2-\rho-\chi\le\rho\chi$, i.e.\ $2\xi+2\eta-\xi\eta\le1$. With $w=\sqrt{\xi\eta}$, AM--GM gives $4w-w^2\le1$, so $w\le2-\sqrt3$ and
\[
\frac{X+Y}{2q^2}=\frac{1+w^2}{2}\le\frac{1+(2-\sqrt3)^2}{2}=4-2\sqrt3=\kappa ,
\]
which is \eqref{eq:B33}.

\emph{Step 4: conclusion.} Combining \eqref{eq:B30}, \eqref{eq:B32}, \eqref{eq:B33} and $\beta\ge\frac{11}{25}q$,
\begin{equation}\label{eq:B34}
T\ \le\ \kappa q^2-\frac{11}{25}q+2 .
\end{equation}
By Lemma~\ref{lem:Hlower}, $H(2q)\ge\kappa q^2-\gamma q-\frac14$, so \eqref{eq:B34} gives $T\le H(2q)$ as soon as
\begin{equation}\label{eq:B35}
\Bigl(\frac{11}{25}-\gamma\Bigr)q\ \ge\ \frac94 .
\end{equation}
Now $\frac{11}{25}-\gamma=\frac1{\sqrt3}-\frac{14}{25}$, so \eqref{eq:B35} at $q=130$ reads $\frac{130}{\sqrt3}-\frac{14\cdot130}{25}\ge\frac94$, i.e.\ $\frac{130}{\sqrt3}\ge\frac{1501}{20}$, whose square is the integer comparison
\[
6{,}760{,}000\ \ge\ 6{,}759{,}003 .
\]
This holds, and the left-hand side of \eqref{eq:B35} is increasing in $q$ (the coefficient $\frac1{\sqrt3}-\frac{14}{25}=0.01735\ldots$ is positive), so \eqref{eq:B35} holds for every $q\ge130$.
\end{proof}

\begin{remark}\label{rem:integrality}
The margin in \eqref{eq:B35} is $0.01735\ldots$ per unit of $q$ against a fixed deficit of $9/4$, which is why the threshold sits at $q=130$ and not lower.  The particular continuous relaxation used before the integrality step is not sharp at $n=18$ and $n=20$: it gives only $t\le43$ and $t\le53$, against the exact values $42$ and $51$.  The integrality step \eqref{eq:B29} --- and, in the finite range, the integer cut envelope of Section~\ref{sec:ladder} --- closes the remaining one- and two-unit gaps.
\end{remark}

\subsection{Proof of the upper bound for \texorpdfstring{$q\ge130$}{q >= 130}}

\begin{theorem}[Large even orders]\label{thm:q130}
$t(2q)\le H(2q)$ for every $q\ge130$.
\end{theorem}

\begin{proof}
Suppose $t(2q)\ge H(2q)+1$. By Lemma~\ref{lem:lines} there is a configuration with $\min(B,W)\ge H(2q)+1$, and by Lemma~\ref{lem:chamber} we may assume it lies in the canonical chamber. Since $q\ge130\ge64$, Corollary~\ref{cor:escape} places its normalized profile in $N$ and its counts in the aggregate core \eqref{eq:core}. Lemma~\ref{lem:cuts} gives $\min(B,W)\le T$, and Propositions~\ref{prop:caseA} and~\ref{prop:caseB} --- which between them cover both orderings of $X$ and $Y$ --- give $T\le H(2q)$. Hence $\min(B,W)\le H(2q)$, a contradiction.
\end{proof}

Explicitly, a canonical profile falls into exactly one of three branches: outside $N$, where Theorem~\ref{thm:cert1} and Lemma~\ref{lem:escape} give $\min(B,W)<H(2q)$; inside $N$ but outside the aggregate core, where Theorem~\ref{thm:cert2} and Lemma~\ref{lem:escape} give the same; and inside the core, where Section~\ref{sec:local} gives $\min(B,W)\le H(2q)$.

\section{The finite range \texorpdfstring{$q\le129$}{q <= 129}}\label{sec:ladder}

\subsection{The certified cut library}

The support cuts used in Section~\ref{sec:local} belong to a library of
$760$.  A cut is a pair $(\theta,\lambda)$ with
$\theta\in\Q\cap[0,1]$ and $\lambda\in\Q^{42}$ such that, on each of the
$64$ atom columns,
\begin{equation}\label{eq:support}
\theta\,\mathbf 1[\varepsilon=(1,1,1,1)]
+(1-\theta)\,\mathbf 1[\varepsilon=(0,0,0,0)]
\ \le\ (\lambda^{\mathsf T}E)_{i,j,\varepsilon}
\end{equation}
for every block $(i,j)$ and atom $\varepsilon\in\{0,1\}^4$, where $E$ is
the $42\times64$ atom-variable coefficient matrix of the moment equations.
Its columns are indexed by $(i,j,\varepsilon)$.

\begin{lemma}[Support-cut inference]\label{lem:cutsound}
If \eqref{eq:support} holds, then every genuine configuration satisfies
\[
\min(B,W)\le\theta B+(1-\theta)W
\le q^2\lambda^{\mathsf T}b(x),
\]
where $b(x)$ is the moment right-hand side at its outer profile.
\end{lemma}

\begin{proof}
The first inequality holds because the middle expression is a convex
combination of $B$ and $W$.  Multiply the columnwise inequalities
\eqref{eq:support} by the nonnegative atom counts and sum.  The moment
equations replace $Ep$ by $b(x)$.
\end{proof}

The archive separates the library into $76$ legacy polynomials in
\texttt{benders\_\bs{}cuts.json} and $684$ delivered rational duals in
\texttt{new\_\bs{}dual\_\bs{}cuts.json}.  An exact standard-library verifier
reconstructs the $42$ moment rows, checks the $64$ column inequalities and
polynomial identity for every delivered dual, and recovers rational
witnesses for the $76$ legacy cuts by exact simplex.  It verifies all
$760$ cuts.

\subsection{The cut envelope at a fixed order}

For fixed $q$, evaluate each cut $P$ at integer outer counts.  Let $s$ be
the least common multiple of $2$ and of all coefficient denominators of the
library; in the released library $s=12$.  The engine stores the integer
coefficients of $sP$, scales constants by $q^2$ and linear terms by $q$, and
performs every comparison as the exact integer test
$s\,P(x)\le s\,H(2q)+s/2$, so the threshold $H(2q)+\frac12$ is represented
without rounding.  The minimum of the cuts is the \emph{cut envelope}.  The
engine explores boxes $[l,u]\subseteq\{0,\ldots,q\}^8$ starting from the full
cube, prunes and tightens by the four chamber and the budget inequalities,
and discharges a box when one cut is certified to stay at most
$H(2q)+\frac12$ on it, by a term-wise interval bound or by an exact corner
maximum; one-point boxes are discharged by direct evaluation.  The following
lemma records why each of these operations is sound.

\begin{lemma}[Box-discharge soundness]\label{lem:boxops}
Throughout the search $0\le l\le u\le q$ coordinatewise, and:
\begin{enumerate}
\item[(i)] \emph{Propagation.}  Each of the five canonical inequalities has
coefficients in $\{0,\pm1\}$.  A box is pruned only when the minimum of a
row over the whole box, evaluated coordinatewise at endpoints, exceeds its
right-hand side; and a bound is tightened only to a value implied by the row
for every point of the box satisfying it.  Hence no integer point satisfying
the inequalities is ever discarded.  Since the root box is the full cube,
every canonical profile is covered; points violating the chamber are covered
as well, which is harmless.
\item[(ii)] \emph{Multi-affinity.}  Every cut polynomial is multi-affine:
its quadratic monomials use only products of two distinct coordinates
(the $22$ admissible pairs), never a square, so it is affine in each
coordinate separately.  The engine rejects any other input shape at parse
time.
\item[(iii)] \emph{Corner maximum.}  A multi-affine function attains its
maximum over a real box at a vertex.  The corner test enumerates all
vertices in the coordinates that occur in the cut and are undetermined in
the box; if none exceeds the threshold, the cut is at most the threshold on
the whole real box, in particular at every integer point.
\item[(iv)] \emph{Interval bound.}  Since $0\le l\le x\le u$, each bilinear
term $c_{ij}x_ix_j$ is at most $c_{ij}u_iu_j$ for $c_{ij}>0$ and at most
$c_{ij}l_il_j$ for $c_{ij}<0$, and each linear term is bounded at its
matching endpoint; the term-wise sum is therefore a valid upper bound for
the cut on the box.
\item[(v)] \emph{Splitting.}  A split at $c=\lfloor(l_i+u_i)/2\rfloor$ into
$x_i\le c$ and $x_i\ge c+1$ partitions the integer points of the parent box.
\item[(vi)] \emph{No overflow.}  All engine arithmetic is signed $64$-bit.
Let $S$ be the maximum over the library of the sum of the absolute values of
all scaled coefficients of one cut; the driver recomputes $S$ before every
run and aborts unless $Sq^2<2^{63}$.  In the released library $S=360$, and
every intermediate quantity of the engine is bounded in absolute value by
$2Sq^2$, which even at the engine's input cap $q\le1000$ is
$7.2\cdot10^{8}\ll2^{63}$; every comparison is therefore exact.
\end{enumerate}
\end{lemma}

\begin{proof}
(i) The minimum of a row over the box is attained coordinatewise at
endpoints because the row is linear; if it exceeds the right-hand side, no
point of the box satisfies the row.  Otherwise, for a coordinate with
coefficient $+1$, every satisfying point obeys
$x_i\le\mathrm{rhs}-(\text{minimum of the remaining terms over the box})$,
which is an integer since all coefficients are $\pm1$; tightening $u_i$ to
it (and $l_i$ symmetrically for coefficient $-1$) keeps every satisfying
point.  (ii) is structural in the stored cut files and is reverified by the
audits.  (iii) Fixing all but one coordinate, a multi-affine function is
affine, hence maximized at an endpoint; induct on the dimension.
(iv) All coordinates are nonnegative because $l$ only ever increases from
$0$, so the displayed per-term extrema are exact.  (v) is immediate.
(vi) For $x\in[0,q]^8$ every term of $sP$ is bounded by its absolute
coefficient times $q^2$, so partial sums are bounded by $Sq^2$; the corner
test's incremental updates combine at most two such sums.  The stated
values of $s$ and $S$ are recomputed by the release audit.
\end{proof}

\begin{lemma}[Integer-box inference]\label{lem:boxsound}
Suppose such a search ends with no survivor and every nonempty terminal box
has a cut bounded above by $H(2q)+\frac12$ throughout the box.  Then
$t(2q)\le H(2q)$.
\end{lemma}

\begin{proof}
By Lemma~\ref{lem:boxops}(v) every split partitions the integer points of
its parent, and by Lemma~\ref{lem:boxops}(i) pruning discards only boxes
containing no canonical integer profile; so induction on the tree shows
that every canonical integer profile belongs to a discharged terminal box.
By Lemma~\ref{lem:boxops}(iii), (iv) --- or by direct evaluation for
one-point boxes --- the discharging cut is at most $H(2q)+\frac12$ at that
profile, and Lemma~\ref{lem:cutsound} gives
$\min(B,W)\le H(2q)+\frac12$ there.  Both $\min(B,W)$ and $H(2q)$ are
integers, hence $\min(B,W)\le H(2q)$.
\end{proof}

\begin{proposition}\label{prop:ladder}
For
\[
q\in\{3,5\}\cup\{10,11,\ldots,129\},
\]
the cut envelope is at most $H(2q)+\frac12$ throughout the canonical integer
domain.  Consequently $t(2q)=H(2q)$ for these orders.
\end{proposition}

\begin{proof}
We ran the released C++17 exact box engine on all $122$ listed orders.  Every
search exhausted its stack with no survivor, visiting $13{,}379{,}092$ nodes
in total; every discharge used one of the $76$ legacy cuts, and the fallback
sweep of the full library at one-point boxes was never invoked.  The release verifier recompiles that engine, repeats every search,
recomputes $H(2q)$ and compares every proof-relevant output field; thus it
re-establishes the domain coverage required by Lemma~\ref{lem:boxsound}.

The archived per-order JSON files bind the cut inputs and record witnesses,
status fields and aggregate statistics, but are compact run summaries, not
node-by-node tree certificates.  Their checker audits those data and the
full-tree accounting without claiming geometric coverage.  As a separate
implementation check, the Python engine produced the original summaries
through $q=60$ and was rerun for every $61\le q\le129$; its non-runtime
envelope outputs agree exactly with the C++ outputs.  The two engines
implement the same box algorithm, so this agreement is corroborative rather
than a conceptually independent proof method.  Lemma~\ref{lem:boxsound}
gives the upper bound, and Proposition~\ref{prop:lower} gives equality.
\end{proof}

The orders $q=1,2$ are elementary.  The four remaining envelope exceptions
$q=4,6,7,9$ are direct finite searches, while $q=8$ is covered by the
union-domain verification of $n=16$.  They are proved together with the odd
searches in Section~\ref{sec:odd}.

\section{The exhaustive finite searches}\label{sec:odd}

On an odd torus the map $\varphi$ of Lemma~\ref{lem:inc} is bijective, so
label parity carries no geometric information.  We now give the finite
arguments for $t(15)$ and $t(17)$ and for the five even orders not covered by
Proposition~\ref{prop:ladder}.

\subsection{The common search reduction}\label{ssec:search}

For odd $n$, write
\[
(s_0,s_1,s_2,s_3)=(|R|,|C|,|D|,|A|),\qquad S=\sum_i s_i.
\]
If a target $M$ is feasible, Lemmas~\ref{lem:budget},
\ref{lem:inc} and~\ref{lem:triple} imply
\begin{align}
S&\le2n,\label{eq:oddprofile1}\\
s_is_j&\ge M,\qquad (n-s_i)(n-s_j)\ge M
&& (i<j),\label{eq:oddprofile2}\\
n^2-n(x+y+z)+xy+xz+yz&\ge2M
&&\text{for each triple }x,y,z.\label{eq:oddprofile3}
\end{align}
These necessary conditions give a finite profile list.  We quotient only by
explicit affine symmetries: transpose, $(r,c)\mapsto(r,-c)$, and, for odd
$n$, $(r,c)\mapsto(r+c,r-c)$.  No classification of the full affine
symmetry group is used.

For each surviving profile, fix two line families, say $R,C$, up to a
subgroup of their affine stabilizer.  For an enumerated set $A$, define
\[
\begin{aligned}
p_d&=\#\{(r,c):r\in R,\ c\in C,\ r-c=d,\ r+c\in A\},\\
q_d&=\#\{(r,c):r\notin R,\ c\notin C,\ r-c=d,\ r+c\notin A\},\qquad
Q_0=\sum_d q_d .
\end{aligned}
\]
For any $D$ of its prescribed size,
\begin{equation}\label{eq:completion}
|\Bl|=\sum_{d\in D}p_d,\qquad
|\Wh|=Q_0-\sum_{d\in D}q_d.
\end{equation}
Thus the last family is an exact fixed-cardinality, two-constraint subset
problem.

\begin{lemma}[Exhaustive-search inference]\label{lem:searchsound}
Assume that the profile filter uses only necessary conditions, the fixed-pair
list meets every orbit of a subgroup preserving the searched profile domain,
every admissible third-family set is enumerated, and
\eqref{eq:completion} is decided exactly for the last family.  If every
listed profile is infeasible, then the target $M$ is infeasible.
\end{lemma}

\begin{proof}
Any feasible line colouring has a profile in the filtered list.  A permitted
affine symmetry sends its fixed pair to an enumerated representative and
transports the remaining line sets within the searched domain.  Its third
set occurs in the enumeration, and its last set satisfies
\eqref{eq:completion}; an exact decision procedure must therefore find it.
\end{proof}

\subsection{The value \texorpdfstring{$t(15)$}{t(15)}}\label{ssec:t15}

The line sets
\[
\begin{aligned}
R&=\{5,7,8,11,12,14\},&
C&=\{0,3,4,7,9,10,12,13\},\\
D&=\{0,1,4,8,12,13,14\},&
A&=\{1,2,3,9,10,12,14\}
\end{aligned}
\]
have $|\Bl|=20$ and $|\Wh|=22$, proving the lower bound.

At target $M=21$, conditions \eqref{eq:oddprofile1}--\eqref{eq:oddprofile3}
leave $247$ canonical profiles.  Their counts in the strata
$S=20,\ldots,30$ are
\[
1,\,1,\,4,\,6,\,13,\,18,\,29,\,37,\,47,\,52,\,39.
\]
For each profile, the fixed-pair subgroup is generated by independent
translations, common multiplication by a unit, transpose when the fixed
sizes agree, and $(R,C)\mapsto(R,-C)$ only when the two completion sizes
agree.  The last restriction ensures that the reflection does not exchange
two distinct ordered profiles.  A separate audit regenerates a canonical
identifier for every orbit and checks all $43$ fixed-pair orbit types by
Burnside's lemma.

The first enumerator decides \eqref{eq:completion} by depth-first search,
using necessary cardinality bounds on each suffix.  An independent
enumerator chooses a different fixed pair on $19$ profiles and uses an exact
$7+8$ meet-in-the-middle Pareto table.  Their respective case totals are
\[
6{,}074{,}753{,}568\qquad\text{and}\qquad6{,}241{,}793{,}402.
\]
Both return \texttt{UNSAT} on all $247$ profiles; the complete second result
file is checked for exact profile coverage and unique \texttt{UNSAT} records.

\begin{theorem}\label{thm:t15}
$t(15)=20$.
\end{theorem}

\begin{proof}
The displayed line sets give $t(15)\ge20$.  The profile generation,
orbit audit, complete third-family enumeration and the two exact
last-family decisions satisfy Lemma~\ref{lem:searchsound} at target $21$.
Hence $t(15)\le20$.
\end{proof}

\subsection{A union-domain verification at \texorpdfstring{$n=16$}{n=16}}
\label{ssec:t16}

Although Theorem~\ref{thm:main} already gives $t(16)=32$, a direct search is
useful both as an independent small-order proof and as the integration of the
earlier $n=16$ computation.  Let $E,O\subset\Z/16\Z$ be the even and odd
residue classes and use profiles
\[
(r,c,d_0,d_1,a_0,a_1)
=(|R|,|C|,|D\cap E|,|D\cap O|,|A\cap E|,|A\cap O|).
\]
The budget, pair bounds and triple identity leave exactly $1{,}898$ ordered
profiles at target $33$.

The proof search groups these profiles only by $(r,c)$, after transposing so
that $r\le c$, and searches the \emph{union} of all associated completion
profiles.  This union is invariant under simultaneous parity exchange and
under exchanging $D$ with $A$.  Hence independent translations, common
multiplication by an odd unit, $(R,C)\mapsto(R,-C)$, and transpose when
$r=c$ act within the searched domain; no parity-orientation transfer is
needed.  Fixed-pair representatives are selected by a lexicographic
canonical-form test on translation necklaces.

The enumerator scans every admissible $A$-mask.  Only the two necessary
row--column--antidiagonal triple inequalities are used as a scalar prefilter;
for every survivor it scans every admissible $D$-mask literally.  The exact
totals are
\[
\begin{array}{@{}lr@{}}
\text{fixed-pair representatives}&159{,}551\\
A\text{-masks tested}&3{,}226{,}530{,}570\\
A\text{-masks surviving the scalar filter}&2{,}951\\
\text{diagonal-cardinality cases}&29{,}768\\
D\text{-masks tested}&38{,}830{,}322 .
\end{array}
\]
None is feasible.  A separately written Python audit regenerates the profile
domain, proves its closure under the three domain symmetries, computes the
seven pair-orbit counts by Burnside's lemma, regenerates every representative
list and hash, and checks the certificate totals.  It prints
\texttt{UNION\_\bs{}AUDIT\_\bs{}OK}.  The older $677$-job programs, which
share their search chassis but use different completion kernels, reproduce
the same upper bound as a secondary cross-check.

\begin{proposition}\label{prop:t16sweep}
The union-domain search certifies that no $33+33$ placement exists on the
$16\times16$ torus.
\end{proposition}

\begin{proof}
The $1{,}898$ profiles contain every feasible colouring by the necessary
conditions.  The union domain makes every quotient symmetry profile
preserving, and the audited representative list meets every fixed-pair
orbit.  All third-family masks and then all admissible last-family masks are
tested, so Lemma~\ref{lem:searchsound} applies.
\end{proof}

\subsection{The remaining orders}\label{ssec:remaining}

\begin{lemma}[Parity-refined even sweeps]\label{lem:evensweep}
Let \(n=2h\), and suppose a line colouring has target \(M\).  With
\(P_\varepsilon=\{x\in\Zn:x\equiv\varepsilon\pmod2\}\), give it the profile
\[
(r,c,d_0,d_1,a_0,a_1)
=\bigl(|R|,|C|,|D\cap P_0|,|D\cap P_1|,
       |A\cap P_0|,|A\cap P_1|\bigr).
\]
Put \(s_R=r\), \(s_C=c\), \(s_D=d_0+d_1\), \(s_A=a_0+a_1\).  For distinct
families \(X,Y\), define
\[
\begin{aligned}
J_{DA}&=2(d_0a_0+d_1a_1),\\
\overline J_{DA}
&=2\bigl((h-d_0)(h-a_0)+(h-d_1)(h-a_1)\bigr),\\
J_{XY}&=s_Xs_Y,\qquad
\overline J_{XY}=(n-s_X)(n-s_Y)
\quad\bigl(\{X,Y\}\ne\{D,A\}\bigr).
\end{aligned}
\]
After possibly complementing all four line sets, every feasible profile
satisfies
\begin{equation}\label{eq:evenfilter}
\sum_X s_X\le2n,\qquad
J_{XY},\overline J_{XY}\ge M,\qquad
n^2-n(s_X+s_Y+s_Z)+J_{XY}+J_{XZ}+J_{YZ}\ge2M
\end{equation}
for every pair and every triple of distinct families.

It is enough to retain one profile from each orbit generated by exchanging
\(R,C\), exchanging \(D,A\), simultaneously swapping \(d_0,d_1\) and
\(a_0,a_1\), and, on the boundary \(\sum_X s_X=2n\), complementation.  For
each retained profile, choose an \(R,C\) and \(D,A\) orientation and search
both simultaneous parity orientations unless they coincide.  Fixing \(R,C\),
one may quotient by independent translations and common multiplication by a
unit, as well as transpose when \(r=c\) and \((R,C)\mapsto(R,-C)\) when
\((d_0,d_1)=(a_0,a_1)\).  Exhausting the resulting pair representatives,
all \(A\) with the prescribed parity sizes, and the exact parity-constrained
\(D\)-completion problem exhausts the target \(M\).
\end{lemma}

\begin{proof}
The first inequality in \eqref{eq:evenfilter} is Lemma~\ref{lem:budget}.
Lemma~\ref{lem:inc} gives \(J_{XY}\) and \(\overline J_{XY}\) as the black
and white pair-intersection counts; the exceptional formulas for \(D,A\)
use its two-or-zero fibres.  Since \(\Bl\) and \(\Wh\) lie in the
corresponding intersections, the pair inequalities follow.  Applying
Lemma~\ref{lem:triple} to each triple gives the last inequality.

Transpose, \((r,c)\mapsto(r,-c)\), and translation by one row induce the
three profile exchanges above; complementation swaps the armies.  An
independent translation of \(R,C\) shifts the \(D,A\) labels by amounts of
the same parity, so it either preserves both parity orientations or
exchanges both.  Every unit modulo \(2h\) is odd and preserves parity.
Transpose and the relative-sign map preserve a job under precisely the
stated size conditions.  Thus the union of the paired orientations is
stable under the fixed-pair quotient and meets every feasible pair orbit.
For each representative, the third-family enumeration uses only necessary
suffix bounds.  Every surviving \(A\) is passed to the exact subset problem
\eqref{eq:completion}, now with the two constraints
\(|D\cap P_\varepsilon|=d_\varepsilon\).  Lemma~\ref{lem:searchsound} then
gives the asserted coverage.
\end{proof}

The same profile--orbit--completion programme was implemented for arbitrary
\(n\le18\), using Lemma~\ref{lem:evensweep} at even orders.  Its independent
audit regenerates every worklist, the Burnside orbit counts and the complete
candidate-space arithmetic.  The upper campaigns needed here are:
\begin{center}
\begin{tabular}{@{}rrrr@{}}
\toprule
$n$ & target & jobs & third-family cases\\
\midrule
$8$  & $9$  & $6$   & $720$\\
$12$ & $19$ & $100$ & $5{,}111{,}848$\\
$14$ & $26$ & $76$  & $70{,}259{,}455$\\
$17$ & $29$ & $316$ & $157{,}158{,}165{,}660$\\
$18$ & $43$ & $259$ & $228{,}644{,}005{,}008$\\
\bottomrule
\end{tabular}
\end{center}
Every job is \texttt{UNSAT}.  For $n=17$, a second implementation with a
different quotient and enumeration order also returns \texttt{UNSAT} on all
$316$ jobs, after $158{,}802{,}287{,}408$ logical third-family cases.

At $n=17$ the line sets
\[
\begin{aligned}
R&=\{0,1,2,4,5,7,8,9\},&
C&=\{0,1,2,4,5,7,8,9,13\},\\
D&=\{2,3,4,7,10,13,14,15\},&
A&=\{2,5,6,7,9,11,12,13,16\}
\end{aligned}
\]
have $|\Bl|=|\Wh|=28$.

\begin{theorem}\label{thm:t17}
$t(17)=28$.
\end{theorem}

\begin{proof}
The displayed sets give the lower bound.  The audited $316$-job campaign
meets the hypotheses of Lemma~\ref{lem:searchsound} at target $29$, giving
the upper bound.
\end{proof}

\begin{proposition}\label{prop:small-even}
For $q\in\{1,2,4,6,7,8,9\}$, one has $t(2q)=H(2q)$.
\end{proposition}

\begin{proof}
For $q=1$, any two distinct cells of the $2\times2$ torus share a row, a
column or a diagonal, so no black--white pair coexists and
$t(2)=0=H(2)$.  For $q=2$, suppose $\min(B,W)\ge3$ for some line colouring
of the $4\times4$ torus.  For each of the pairs $(R,C)$, $(R,D)$, $(C,D)$
and $(R,A)$ --- none of which is the pair $\{D,A\}$ --- Lemma~\ref{lem:inc}
matches the label pairs bijectively with cells, so with $x,y$ the two black
sizes, $xy\ge3$ and $(4-x)(4-y)\ge3$; the only integer solutions have
$x+y=4$.  The first three pairs force $|R|=|C|=|D|=2$, and the fourth then
gives $|A|=2$.  The triple identity \eqref{eq:triple} for $R,C,D$ now reads
$U+U'=16-4\cdot6+3\cdot4=4<6\le|\Bl|+|\Wh|$, a contradiction.  Hence
$t(4)\le2=H(4)$.  For $q=4,6,7,9$, the first,
second, third and fifth rows of the campaign table exclude
$H(2q)+1$ queens of each colour.  Proposition~\ref{prop:t16sweep} does the
same for $q=8$.  Proposition~\ref{prop:lower} supplies every lower bound.
\end{proof}

\begin{proof}[Proof of Theorems~\ref{thm:main} and~\ref{thm:oddmain}]
Proposition~\ref{prop:lower} gives $t(2q)\ge H(2q)$.  The reverse inequality
is Theorem~\ref{thm:q130} for $q\ge130$, Proposition~\ref{prop:ladder} for
$q\in\{3,5\}\cup\{10,\ldots,129\}$, and
Proposition~\ref{prop:small-even} for the remaining $q$.  Theorems
\ref{thm:t15} and~\ref{thm:t17} give Theorem~\ref{thm:oddmain}.
\end{proof}

\begin{proof}[Proof of the corollaries]
Corollary~\ref{cor:asym} follows from Theorem~\ref{thm:main} and
Lemmas~\ref{lem:Hupper} and~\ref{lem:Hlower}.  The displayed values follow
by evaluating the finite maximum \eqref{eq:Hdef}.
\end{proof}

The odd construction below explains some of the sensitivity to divisors, but
does not suggest an odd analogue of Theorem~\ref{thm:main}.

\begin{proposition}[Divisor lifting]\label{prop:lift}
If $d\mid n$ then $t(n)\ge(n/d)^2\,t(d)$.
\end{proposition}

\begin{proof}
Let $(R,C,D,A)$ be an optimal configuration on the $d\times d$ torus and let $\pi:\Zn\to\Z/d\Z$ be reduction. Take $(\pi^{-1}R,\pi^{-1}C,\pi^{-1}D,\pi^{-1}A)$. A cell $(r,c)$ of the $n$-torus is black-homogeneous iff $(\pi r,\pi c)$ is, since $\pi(r-c)=\pi r-\pi c$ and $\pi(r+c)=\pi r+\pi c$; and each cell of the $d$-torus has exactly $(n/d)^2$ preimages. So both armies scale by $(n/d)^2$, and Lemma~\ref{lem:lines} applies.
\end{proof}

For example, $t(39)\ge144$, $t(45)\ge180$ and $t(51)\ge252$.  The lift from
proper divisors is not generally sharp:
\[
\max_{\substack{d\mid9\\d<9}}(9/d)^2t(d)=0<7=t(9),\qquad
\max_{\substack{d\mid15\\d<15}}(15/d)^2t(d)=18<20=t(15).
\]
The odd asymptotics remain open; the upper bound $t(n)\le n^2/8$ is
\cite[Thm.~1.4]{CDHS24}.

\section{Use of AI tools}\label{sec:ai}

Large-language-model tools assisted with conjecture generation, drafting,
code generation and source review.  The author checked the final mathematical
arguments and proof-critical outputs and takes responsibility for the paper.
Agreement among AI-assisted implementations is treated as an error-detection
measure, not as external scholarly reproduction.

\section*{Acknowledgements}

The work of Clinch, Drescher, Huynh and Saffidine provided the decisive plaid
starting point and the previous asymptotic bounds.  I also thank the
contributors who have maintained OEIS A279405 and the reviewer whose
incidence correction and union-domain verification materially strengthened
this paper.  A second independent referee replayed the exact certificates and
the complete finite ladder, prompted the formal dual-leaf and box-discharge
inference lemmas, and caught a packaging defect in a distributed copy of the
ancillary archive.

\appendix

\section{Verification and reproduction}\label{app:verify}

\sloppy

The complete ancillary archive is bound by \texttt{SHA256SUMS}.  The archive
is packaged without platform metadata files, and the manifest check of the
release verifier is two-sided: it fails on any missing, altered or unlisted
file, ignoring only byte-code caches.  An immutable
public snapshot is available at
\begin{center}
\url{https://github.com/jphme/math-problems/tree/459046273420a61d51bac18f3a260b4317ec2e62/peaceable-queens-even-torus}.
\end{center}
The following command, run from \texttt{anc/}, replays the exact proof
objects, reruns the complete finite cut envelope and audits the frozen
outputs of the larger searches:
\begin{verbatim}
uv run --isolated --python 3.14.6 \
  --with sympy==1.14.0 --with mpmath==1.3.0 \
  python verify_bundle.py
\end{verbatim}
It ends with
\begin{center}
\texttt{PEACEABLE\_\bs{}QUEENS\_\bs{}RELEASE\_\bs{}BUNDLE\_\bs{}OK}.
\end{center}
The recorded environment is CPython~3.14.6 and SymPy~1.14.0; all other
Python checks use only the standard library.  The sources use C++17 or C++20
as specified by their build commands and were rebuilt with Apple
Clang~21.0.0; the union-domain enumerator was also reproduced with GCC~14 and
Clang~17.  No floating-point comparison affects a proof verdict.

\subsection*{A.1 Exact proof objects}

The files \texttt{even\_\bs{}c527\_\bs{}sym\_\bs{}mc.json.gz} and
\texttt{even\_\bs{}local\_\bs{}core\_\bs{}mc.json.gz} contain the two branch
trees.  The delivered checker and a separately written verifier both
reconstruct every split, moment row, McCormick inequality and rational leaf
dual.  The latter also derives the $42$ moment rows from torus incidence
fibres and tests them on explicit colourings.

The files \texttt{new\_\bs{}dual\_\bs{}cuts.json} and
\texttt{benders\_\bs{}cuts.json} together contain the $760$ support cuts.
\texttt{verify\_\bs{}support\_\bs{}duals.py} verifies all of them exactly and
recovers the $76$ legacy duals independently.  The two cuts used in
Section~\ref{sec:local} are also checked by
\texttt{verify\_\bs{}even\_\bs{}finite\_\bs{}algebra.py}; the independent
prose checker verifies the displayed algebra and exact constants.

\subsection*{A.2 Finite computations}

The directories \texttt{records\_\bs{}q003\_\bs{}q020/} and
\texttt{records\_\bs{}q021\_\bs{}q129/} contain compact summaries of the
$122$ ladder runs.  They do not contain the split nodes or terminal boxes.
\texttt{check\_\bs{}envelope\_\bs{}records.py} recomputes $H$ and checks the
order set, schemas, source and input hashes, thresholds, witnesses and
aggregate node/leaf accounting.  Full coverage is re-established by
\begin{verbatim}
python replay_finite_envelope.py
\end{verbatim}
This compiles \texttt{box\_\bs{}engine.cpp} with strict C++17 flags and
repeats all $122$ searches.  The separately written Python implementation
was also rerun for $q=61,\ldots,129$; its exact agreement with the C++ output
is frozen in
\texttt{finite\_\bs{}envelope\_\bs{}python\_\bs{}crosscheck\_\bs{}q061\_\bs{}q129.json}.
The ledger checker audits that frozen cross-check, while
\texttt{crosscheck\_\bs{}finite\_\bs{}backends.py} recreates it.  These
programs share the same exact integer-box algorithm.

\begin{table}[ht]
\centering\small
\begin{tabular}{@{}
>{\raggedright\arraybackslash}p{0.20\textwidth}
>{\raggedright\arraybackslash}p{0.32\textwidth}
>{\raggedright\arraybackslash}p{0.37\textwidth}@{}}
\toprule
Range & Full computation & Action in the one-command verification\\
\midrule
$q=3,5,10,\ldots,129$ &
\texttt{box\_\bs{}engine.cpp} &
full exact rerun and comparison with every frozen summary\\
$n=15$ &
\texttt{peace15\_\bs{}solver.cpp}, \texttt{profile\_\bs{}enum.cpp} &
profile, orbit, output-coverage and witness audits\\
$n=16$ &
\texttt{peace16\_\bs{}union\_\bs{}enum.cpp} &
domain, representative, hash, total and witness audits\\
$n=8,12,14,17,18$ &
\texttt{general\_\bs{}sweep/solver\_\bs{}b.cpp} &
worklist, orbit, output-coverage and witness audits\\
\bottomrule
\end{tabular}
\caption{Verification modes for the finite searches.  The README gives the
full rebuild commands for the computations audited only through frozen
outputs here.}\label{tab:finite-modes}
\end{table}

For $n=15$, the archive includes both enumerator sources, both complete
outputs, the $247$-profile worklist, a Burnside audit and direct witness
checks.  For $n=16$, it includes the union-domain source, certificate and
independent representative-list audit, as well as the older two-kernel
cross-check.  The directory \texttt{general\_\bs{}sweep/} contains the
general-$n$ source, audits and frozen outputs used for $n=8,12,14,17,18$,
together with the second $n=17$ implementation and output.

The rational branch trees and support duals are standalone proof objects.
The ladder summaries and exhaustive-search outputs are audit records: their
logical force is the corresponding short source together with the reductions
proved in Sections~\ref{sec:ladder} and~\ref{sec:odd}.  The audits establish
the necessary profile domains, the orbit and representative coverage, the
candidate-space totals, the consistency and uniqueness of the output records,
and the positive witnesses; they do not certify each individual completion
test of the large enumerations.  Among the proof-critical campaigns, $n=15$
and $n=17$ were decided twice by independently written enumerators, and
$n=16$ additionally by the older two-kernel programs; the $n=8,12,14,18$
sweeps rest on the single audited implementation
\texttt{general\_\bs{}sweep/solver\_\bs{}b.cpp}.  The release command reruns
the complete ladder; the larger board enumerations are reproducible with the
complete build and run commands in \texttt{README.txt}, which also records
the machine, compiler, per-campaign case totals and wall-clock data, and the
SHA-256 digests of the frozen outputs.  Emitting independently checkable
UNSAT certificates for these enumerations, and a second implementation for
the $n=18$ sweep, remain natural hardening steps.

\fussy

\begin{thebibliography}{99}

\bibitem{Ainley77}
S.~Ainley,
\emph{Mathematical Puzzles},
G.~Bell \& Sons, London, 1977.

\bibitem{CDHS24}
K.~Clinch, M.~Drescher, T.~Huynh, and A.~Saffidine,
\emph{Constructions, bounds, and algorithms for peaceable queens},
extended version, arXiv:2406.06974v3 (2024).

\bibitem{CDHS25}
K.~Clinch, M.~Drescher, T.~Huynh, and A.~Saffidine,
\emph{Constructions, bounds, and algorithms for peaceable queens},
in \emph{2025 Proceedings of the Symposium on Algorithm Engineering and
Experiments (ALENEX)}, SIAM, 2025, pp.~160--167,
\href{https://doi.org/10.1137/1.9781611978339.13}
{doi:10.1137/1.9781611978339.13}.

\bibitem{OEIS279405}
N.~J.~A. Sloane and others,
\emph{The On-Line Encyclopedia of Integer Sequences}, sequence A279405:
\emph{Peaceable coexisting armies of queens on a toroidal board},
\url{https://oeis.org/A279405}, accessed July 27, 2026.

\end{thebibliography}

\end{document}
